\documentclass[11pt,letterpaper]{article}
\def\SuiteDocTitle{Release and Iteration Planning Procedure}
\def\SuiteDocCode{PR-RELEASE}
% Shared professional style: pdfLaTeX compatible, no external font files.
\usepackage[T1]{fontenc}
\usepackage[utf8]{inputenc}
\usepackage{lmodern}
\usepackage{microtype}
\usepackage[margin=0.83in,headheight=16pt,footskip=28pt]{geometry}
\usepackage[table]{xcolor}
\definecolor{navy}{HTML}{12304A}
\definecolor{teal}{HTML}{087F83}
\definecolor{gold}{HTML}{B78633}
\definecolor{blueback}{HTML}{F0F5F8}
\definecolor{tealback}{HTML}{EEF8F6}
\definecolor{amberback}{HTML}{FFF7E8}
\definecolor{grayback}{HTML}{F6F7F8}
\usepackage{array,tabularx,longtable,booktabs}
\usepackage{enumitem}
\usepackage{amssymb}
\usepackage{titlesec}
\usepackage{fancyhdr,lastpage}
\usepackage{tcolorbox}
\tcbuselibrary{breakable}
\usepackage{pgffor}
\usepackage{hyperref}
\hypersetup{colorlinks=true,linkcolor=teal,urlcolor=teal,pdfauthor={Engineering Learning Series},pdftitle={\SuiteDocTitle}}
\setlength{\parindent}{0pt}
\setlength{\parskip}{5.5pt}
\setlist{topsep=3pt}
\setcounter{tocdepth}{2}
\titleformat{\section}{\Large\bfseries\color{navy}}{\thesection}{0.7em}{}
\titleformat{\subsection}{\large\bfseries\color{teal}}{\thesubsection}{0.6em}{}
\titleformat{\subsubsection}{\normalsize\bfseries\color{navy}}{\thesubsubsection}{0.6em}{}
\pagestyle{fancy}
\fancyhf{}
\fancyhead[L]{\scriptsize\color{navy}AGILE REQUIREMENTS ENGINEERING}
\fancyhead[R]{\scriptsize\color{teal}\SuiteDocCode}
\fancyfoot[L]{\scriptsize Independent learning and implementation materials}
\fancyfoot[R]{\scriptsize Page \thepage\ of \pageref{LastPage}}
\renewcommand{\headrulewidth}{0.35pt}
\renewcommand{\footrulewidth}{0pt}
\newcommand{\SuiteHeader}[2]{%
{\color{teal}\small\bfseries ENGINEERING DOCUMENTATION SUITE \hfill #2}\par\vspace{8pt}
{\color{navy}\LARGE\bfseries\SuiteDocTitle}\par\vspace{5pt}
{\color{black!68}\normalsize #1}\par\vspace{8pt}
{\color{gold}\rule{\linewidth}{2pt}}\par\vspace{8pt}}
\newcommand{\WritingLines}[1]{\par\foreach \n in {1,...,#1}{\noindent\rule{\linewidth}{0.35pt}\par\vspace{0.28cm}}}
\newcommand{\FreeSpace}[1]{\begin{tcolorbox}[colback=white,colframe=black!25,height=#1,boxrule=0.5pt]\end{tcolorbox}}
\newcommand{\Field}[1]{\noindent\textbf{#1}\quad\rule{0.42\linewidth}{0.35pt}\par\vspace{0.3cm}}
\newcommand{\CheckItem}[1]{\noindent$\square$\quad #1\par\vspace{2pt}}
\newcommand{\SmallNote}[1]{{\footnotesize\color{black!62}#1}\par}
\newcommand{\Guidance}[1]{\begin{tcolorbox}[colback=blueback,colframe=teal,breakable,title=Guidance]#1\end{tcolorbox}}
\sloppy

\begin{document}
\SuiteHeader{Operating procedure | responsibilities, decisions, controls and evidence}{Procedure}
\tableofcontents\medskip\noindent\rule{\linewidth}{0.4pt}\medskip
\section{Purpose}
Create an evidence-based, revisable forecast and inspect delivered user outcomes.

\section{Trigger and entry conditions}
\begin{itemize}[leftmargin=*,itemsep=2pt]
\item Release objective, available team, prioritized backlog and relevant historic capacity signals are known.
\end{itemize}
\section{Roles and decision rights}
\begin{longtable}{|p{0.28\linewidth}|p{0.60\linewidth}|}\hline
\textbf{Role} & \textbf{Responsibility} \\ \hline\endfirsthead\hline
\textbf{Role} & \textbf{Responsibility} \\ \hline\endhead
Product representative / authorized stakeholder & Owns value, ordering and intended acceptance \\ \hline
Delivery team & Estimates, slices, implements and verifies technical behavior \\ \hline
Tester / quality specialist & Helps design objective examples and regression evidence \\ \hline
Security / risk representative & Advises on security-sensitive constraints and exceptions \\ \hline
Facilitator & Keeps assumptions, decisions and follow-up visible \\ \hline
\end{longtable}
\section{Operating steps}
\begin{enumerate}[leftmargin=*,itemsep=3pt]
\item Negotiate release goal, date or scope boundary and key operational risks.
\item Order stories based on user value, risk and external dependencies.
\item Forecast using team-local observations and present at least one alternative scope scenario.
\item For each iteration, confirm capacity, sprint/iteration goal, test readiness and technical tasks.
\item Inspect completed and accepted outcomes regularly; monitor blocked work and scope movement.
\item At review, demonstrate increments and obtain product feedback.
\item Update plans, constraints and acceptance evidence, recording material decisions.
\end{enumerate}
\section{Exit criteria}
\begin{itemize}[leftmargin=*,itemsep=2pt]
\item Accepted increments, updated forecast, known risks and next action are visible.
\end{itemize}
\section{Exception / escalation}
For unclear authority, conflicting business rules, urgent production impact or unacceptable security risk, stop implicit acceptance. Identify the decision owner, time-box the investigation, document the temporary mitigation and escalate under the organization’s established governance. A process exception must not silently override a binding control.

\section{Evidence and quality checks}
\begin{itemize}[leftmargin=*,itemsep=2pt]
\item Retain an item ID, source, owner, relevant timestamps, decisions and tests.
\item Cross-link requirements, stories, control obligations and implementation evidence where risk warrants it.
\item Sample decisions after a review cycle and update the procedure if recurring defects or bottlenecks appear.
\end{itemize}
\section{Basis and extension}
The collaboration, progressive elaboration, estimation and acceptance concepts follow Cohn (2004), especially Chapters 4–11 and 14–16. The governance and control structure is an original modern operational extension.


\end{document}
